\documentclass[../principal.tex]{subfiles}
\graphicspath{{\subfix{../imagenes/}}}

\begin{document}

\chapter{Herramientas}

Creemos que es importante compartir las herramientas que hacen posible nuestra práctica artística, que mucahs de ellas son elegidas por sus bajos costos, la robustez y el cariño de las comunidades que documentan y desarrollan, y que permiten que nuestros computadores sean nuestro taller móvil.

\section{LaTeX}

Este libro está programado en LaTeX. Esta herramienta tiene una larga tradición de uso en ciencias exactas y aplicadas, debido a su capacidad de citar bibliografías de forma ordenada, de crear figuras citables, y de mantener en orden proyectos muy complejos a nivel de edición.

Para este libro estamos usando adicionalmente las subherramientas Tikz y Circuitikz, que permiten crear gráficos y diagramas de circuitos electrónicos.

\section{Git}

Somos gente entusiasta por el control de versiones, por lo que este libro está hecho usando Git, una herramienta de control de versiones que permite tener no solamente una versión actual del libro, sino que también una bitácora de cada cambio.

\section{GitHub}

GitHub es una empresa que permite el almacenamiento en la nube de repositorios con control de versiones de Git.

Además estamos usando la herramienta GitHub Actions, que permite automatizar el proceso de compilación del código fuente de este libro, y su exportación a PDF, y ser subido a un enlace permanente donde cualquier persona puede consultar la versión más reciente de este documento.

\section{KiCad}

Kicad es un software gratuito y de fuente abierta para diseño y fabricación de circuitos electrónicos.

Es una colección de herramientas, que permiten dos grandes 

\section{C++}

C++ es un lenguaje de programación, que hereda la complejidad del lenguaje C, y le agrega la capacidad de crear clases, y con eso, un paradigma llamado programación orientada a objetos. 

C++ es un lenguaje que en nuestra práctica artística cotidiana no usamos, ya que nos encontramos usando scripting en otros lenguajes más modernos, como JavaScript para web, o Python para automatización en general.

\newpage

\end{document}

